\chapter{培养体系}
\section{培养方案}
\SourceParagraph{南赫学院的培养方案整体来说还是比较繁忙的。基础要求同其他学院一样，各种思政课、体育课、通识课都不少，但是大多数在前两年都能结束。}
\SourceParagraph{南赫学院的各类课程主要分为以下几类：基础课，概论课，专业核心，地球系统学科交叉模块，自选模块，专业选修。基础课包括各类英语课、编程课、数学课、物理课，主要集中于前两年，但是大三大四各有一门 4 学分的英语课；概论课程包括大气科学概论、环境科学概论等，全部集中于大二；与之相应的专业核心课程，与概论课程一脉相承，包括动力气象、气候学基础、天气学原理等，全部集中于大三；接下来的三个部分，看上去是自选课程，事实上每个模块都有必修课程，如果只想达到毕业要求，能自由选择的只有一门课或两门课，详情请见附录。除此之外，还有各种杂七杂八的课程，比如实验课（大二下暑假），毕业论文，或者各类学科竞赛、科考活动认定的学分。}
\SourceParagraph{附录：南赫学院各模块课程，上色为必修。}
\SourceParagraph{（基于 2023 级本科培养方案，事实上 23 级开课顺序不完全符合培养方案，24、25 级不排除培养方案被进一步更改的可能性）}
\clearpage
\begingroup
\setlength{\paperwidth}{420mm}
\setlength{\paperheight}{297mm}
\setlength{\textwidth}{392mm}
\setlength{\textheight}{269mm}
\setlength{\oddsidemargin}{-18.4mm}
\setlength{\evensidemargin}{-18.4mm}
\setlength{\topmargin}{-18.4mm}
\setlength{\headheight}{0pt}
\setlength{\headsep}{0pt}
\setlength{\footskip}{8mm}
\setlength{\linewidth}{\textwidth}
\setlength{\hsize}{\textwidth}
\setlength{\vsize}{\textheight}
\special{papersize=420mm,297mm}
\thispagestyle{plain}
\vspace*{\fill}
\begin{center}
  \includegraphics[width=0.995\textwidth,height=0.94\textheight,keepaspectratio]{figures/p25\_01.png}
\end{center}
\vspace*{\fill}
\clearpage
\endgroup
\special{papersize=210mm,297mm}
\section{课程介绍}
\subsection{大一上}
\subsubsection{微积分 I}
\SourceListItem{1，目前了解到微积分 I 改为了由周晨老师授课，24 级微积分 II 是他讲的，建议把微积分 II 的食用指南先看一遍。目前了解到 23 级和 24 级的微积分 I 难度相差很大。25 级及以后的微积分 I 和 II 应该都是周晨老师授课了。}
\SourceLabel{2，难喝的微积分 I 授课难度易于第一层次}
\SourceListItem{3，不出预料这将是你上大学后第一门专业课，请认真对待。}
\subsubsection{大学化学}
\SourceLabel{食用指南：}
\SourceListItem{1. 该指南主要针对南赫大学化学课程，并且使用前提是授课老师仍然是肖守军、袁帅和谢然三人或者课程体系并未大改。}
\SourceListItem{2. 在南赫资料分享群有往年大化的资料，可以下载使用。作业自己完成，建议做作业时要有自己的思考。}
\SourceListItem{3. 如果你是在高中学过大学化学，并且不是“半吊子”的那种，理论上你可以不用看这个了（），把自己之前学的无机知识捡回来就行（虽然课程有有机内容，但是那个显然太简单了，把命名背会就行）。只要你不犯大错误，基本没啥问题，。}
\SourceParagraph{我会从课程内容与难度、上课老师与课堂情况、如何尽量学好这门课（即如何尽量拿一个相对高的分数）、推荐书目\&课程这四个方面进行说明：}
\SourceLabel{1. 课程内容与难度：}
\SourceParagraph{大学化学课程主要内容包括结构、化学原理（简单版物化？）、元素和少量有机知识，一学期学完。具体知识可以看课本，前两年使用的课本都是 Chemistry (Steven S.Zumdahl, Susan A. Zumdahl )这本书，你们上课之前会发一本实体书的。}
\SourceParagraph{实际难度感觉中等，不过感觉大多数同学一开始会有点不适应，因为大学化学和高中化学完全不是一回事，由此也不要担心如果高中化学学不好的话，会不会很麻烦。对于这个担心我只能说，真没啥关系，只要好好学也可以的。}
\SourceLabel{2. 上课老师与课堂情况：}
\SourceParagraph{大学化学由三个老师负责：肖守军老师、袁帅老师和谢然老师（26 级开始有所变动）。这三个老师都是化学化工学院的老师。肖守军老师主要负责前半学期的结构部分、袁帅老师主要负责化学原理部分而谢然老师主要负责元素和有机初步部分，三个老师人都很好，问他们问题完全没问题。}
\SourceParagraph{上课是用英语上的，ppt 和作业也是全英的。在英语这个问题上，肖守军老师的口音可能较难适应，而另外两个老师英语的口语听清英文单词应该没啥问题。}
\SourceParagraph{对于 24 级的课程情况，上课到位率还算不错（毕竟老师讲得都还不错，感觉是真能学到东西的），然后分数组成比例是平时分 30\%，期中 30\%-35\%，期末 35\%-40\%（按照情况调整）。一般来说，作业交了并且不是乱写的，平时分就不会怎么扣，基本都是 29-30 分。至于老师说的在 qq 群多问问题可以加平时分，我个人感觉没啥用，可能只是为了营造一种学习氛围吧。}
\SourceParagraph{从 24 级的情况来看，期中和期末难度都还算正常吧，期中平均分接近 80，期末平均分未知，但是我感觉总评的平均应该在 83 左右，该捞的应该也捞了，只有完全不学才会挂科。}
\SourceLabel{3. 如何尽量学好这门课：}
\SourceParagraph{学好这门课的核心还是认真听讲，不会的内容多与任课老师、助教以及同学进行讨论。我个人更加建议与认识的之前学过大学化学的同学讨论讨论，只要不是“半吊子”，他们都会的。这样做还有一个优势是，他们有很多书可以翻，而不局限与课本，可能能在某本书上找到更好的说法，助于理解。}
\SourceParagraph{对于如何拿更高分这个问题，功利的来讲，只要拿满平时分、考好期中和期末就够了（什么废话），对于拿满平时分来说，上面已经提及了，只要作业认真写了以及上课翘课没被发现就不会扣分（认真写了\ensuremath{\neq}一定要全对，重点是要有自己的思考，而不是抄作业，其实抄袭一眼就能看出来）。}
\SourceParagraph{对于期中考试和期末考试，期中考试就只考结构部分内容，而期末考是不考前半学期内容的，主要就是化学原理、元素和少量有机。考试题型有选择题、填空题和解答题，主要就是平常练习的难度，不会太难的。}
\SourceParagraph{按照老师的说法，只要把作业搞明白，有时间刷刷教材上的习题，考试应该都没问题的，确实也是如此。但如果你想拿更高的分数的话，可以找一下大学化学原理 A、B 的往年卷来做一下（我觉得 nju 的化学就是有一个题库的），去年这么做会发现考试是能做到原题的。考试的时候记得带计算器，至于其他的，只能祝你们好运了。}
\SourceLabel{4. 推荐书目\&课程：}
\SourceParagraph{鉴于化学这门课程就是需要大量阅读各类化学书籍的，并且以防各位无法理解英文书\&ppt，下面我推荐一些中文书籍（按照优先程度排序，最前面的最重要）：卞江等《普通化学原理》、宋天佑等《无机化学》、北师大等《无机化学》、张祖德《无机化学》、南京大学《大学化学》，如果平常没有什么时间，阅读前两本即可，后面的书可能有点深入了。（出于尊重，我把 nju 自己编写的《大学化学》放在这里了，实际这本书有一堆问题，建议别看）}
\SourceParagraph{我感觉 b 站上的网课其实没啥能把大学化学这门课给讲明白的，实在要看的话，去搜索宋天佑的大学化学上课视频就行了。}
\SourceParagraph{如果担心英文化学听不懂的话，想再在南大招一门中文化学课旁听，这是可以的，在此我建议旁听现代工程与应用科学学院的大学化学（旁听\ensuremath{\neq}选课），那个课程体系与我们的基本重合（感觉现工院大化=南赫大化中文版），尽量不要去听化学化工学院的大学化学 A（因为你的诉求不就是要听懂化学课程吗？），那门课相对难度更高，并且体系不同，从知识角度，你只能学到一半的内容}
\subsubsection{大学英语}
\SourceParagraph{南赫的听说和读写课是分成 A 班和 B 班的。A 班对应其他专业二层次的内容，B 班是三层次，并且你们应该是不能换老师的（和其他专业同学不一样）。这里尽量写一点。}
\SourceListItem{1，AB 班之间是可以自己换的，可以根据自身水平和红黑榜（Note: 红黑榜详见下面学习建议 - 选课 处）调整。不要觉得自己分到 A 班就很好，分到 B 班就完蛋了。首先，在计算分数的时候，是不折算直接算入的，比如一个 A 班同学是 85 分，一个 B 班同学是 86，那就是 B 班同学更高。而且 B 班（三层次）的同学课程内容更简单，事情更少,因此这个是需要权衡利弊的（虽然是学校分配的）。}
\SourceListItem{2，虽然英语课和理科课程一样，有要考试算分的部分。但是主观分依然占很大部分，一些刷存在感，多互动之类的方法可能有效，但是不保证有效。}
\SourceListItem{3，卷英语的性价比不大，但是要是你有卷绩需求还是要好好学的，毕竟大一英语占八个学分。}
\SourceListItem{4、有的英语课是有课本且按照课本走的，还需要网测，资源分享群里有网测复习的题库，往年有抽到过原题。}
\SourceListItem{5，背考纲单词的性价比很低，接近 4k 单词只考期末的十分，不过话说回来，想学英语的话还是得自己背背单词。}
\SourceListItem{6、推荐一些 ai 工具，如各类大语言模型（deepseek，gpt，gork），翻译模型（deepl），降 ai 率模型（the spear）、同声翻译（有道留学听课堂）。有的要付费有的要魔法，请妥善按需使用。}
\SourceParagraph{（据传大学英语可能会换新教材，如若传言为真，请慎重对待这里的建议）}
\SourceParagraph{ 25级开始用的是新教材，只有数字教材。根据老师的说法比之前难了不少，个人感觉难度主要体现在里面的音频不是为了英语教材专门录的，而是从各种演讲访谈脱口秀等中截取的，所以会有些比较糊的地方或者背景音乐声音把内容覆盖掉的情况。因此不论英语基础如何请认真准备上机考试，上机考试的成绩和复习认真程度的关系高于和英语水平的关系。听说教材script有的老师会给，如果都没有官方的，有些同学AI或者爬虫出来的就特别准，有的则错漏百出，所以可以同学相互间分享下。}
\SourceParagraph{以下部分只针对想要95+的同学，如果不考虑成绩，听英语纯粹就是浪费时间，因为大家学习方法和基础都各异，几乎不可能出现老师讲课的节奏和难度都最适合你的情况，这点相信大家在高中就有所体会。而且即使课上完全不听不互动，绝大多数老师也不会在平时分上为难学生，因为也没有什么人听，人在，pre好好准备即可。}
\SourceParagraph{关于平时分：如果老师是直接点名型的，那么老师点到你的时候好好回答就很出彩了，因为大部分同学站起来都不知道老师在讲哪里，能有一点自己的想法更好。这种情况如果是读写一学期也就被叫个两三次，有个小技巧是期末之前的那次认真回答，因为老师如果没有每次做记录的话到给分的时候如果正好想到你最后一次回答和期末之前两周的认真，之前稍微摸点鱼完全不影响。如果老师是主动回答型的，那么有几次也够了，特别多也没用。}
\SourceParagraph{关于作业，听说作业如果只在WElearn上做，只要记得做就行（这个平台特别离谱，可以刷到满分为止）；如果是在智慧南雍之类上只有一次机会的，如果没有把握可以去找已经交了的同学对答案。}
\SourceParagraph{关于pre：选好队友。个人感觉第一学期听同学pre的唯一价值就是第二个学期可以找些好队友。}
\SourceParagraph{关于考试：听力四级难度，如果四级轻易全对应付期末完全没问题，如果觉得有难度就按照自己的方法练，无非就是泛听听写做题，如果觉得枯燥可以去看国外的综艺访谈，直接去外网看那种b站过不了审的生肉片段，完美解决学习动力问题；阅读大家应该都十分轻松；小猫题无非就是在非第一词义和生词上为难人（25届上学期有whereas 和whereby 的辨析，个人感觉是出的最刁钻的）；词汇题可能是丢分最严重的题，主要考的就是词义，好像也没考偏门的意思，有几题考的是用法和搭配，但大部分就是为了看这个词你有没有背出来。
}
\SourceBlankLine[2]
\subsubsection{思政课（这是一个统称）}
\SourceParagraph{思政课最重要的就是，选个好老师！好老师真的会决定你的思政课成绩！！！因此要做的事情就是看着红黑榜找好老师，如果分配到的老师本身还可以，那就没必要换了，如果是黑榜，那就一定要换！因为这种政治课实际重要性与大学化学差不多，好的老师可以轻松让你 90+，坏的老师可能让你直接不到 80，而且上课也很累。（在绩点计算法中你会知道这点很重要）}
\SourceParagraph{最开始教务员会给你指选一个思政老师，而在退补选的时候是可以换的。南赫学院分配的思政课的老师一般都是红榜，没有特殊情况的话换课需谨慎（如选了其他课与思政课时间冲突等）。退补选换课就是拼手速，比谁手速快和网速快，详细的退补选策略不过多展开。}
\SourceBlankLine[2]
\subsubsection{体育、北欧文化与社会、认识地球系统科学}
\SourceParagraph{体育课是全校所有学生都可以选的，建议先参考红黑榜，再考虑个人兴趣，选课时，可以选择三个体育课，但是中签只能中一个，而且没中签的概率不小。这时就要靠退补选抢课了。}
\SourceParagraph{剩下三门课都是真正的水课，点名记得去之外，记得完成任务。对大气感兴趣的同学，可以听一下认识地球系统科学这门课，这门课上课方式就是老师请各地的专家线上/线下做讲座，感兴趣可以听一下，期末写 3000 字英文论文。可以让各位对大气科研领域有初步了解。}
\SourceBlankLine[2]
\subsection{大一下}
\subsubsection{微积分 II}
\SourceLabel{食用指南：}
\SourceListItem{1，看课程体系是否更改，老师是否更改。如果更改较多，细致性的指导就没有用了。}
\SourceParagraph{教学老师：周晨（中国人，口语听起来 Chinglish，易于理解，课后交流可以用英文）}
\SourceParagraph{教材：jams-stewart-calculus-early- transcendentals-7th （微积分 II 教了 11-16，微分方程没学）}
\SourceLabel{课堂评价：}
\SourceParagraph{周老师因为是中国人的关系，口语不够 native，反而理解方面完全没有问题，词汇简单，还慢悠悠的，如果觉得下面听不懂还会用中文再说一下。但是有点没啥活力和激情。平时作业和资料会发在一个叫“课堂派”的平台。}
\SourceLabel{给分：}
\SourceParagraph{作业 15\%，作业都是课本题，有十分完整的参考答案，但是量有一点点大。}
\SourceParagraph{随堂检测 15\%，也就是我们所说的 quiz，24 届的 quiz 定在每周四的最后一节课（也就是这一周的微积分课中的最后一节），请不要翘这一堂课，当然狠人也可以只上这一堂课。quiz 是对上一周的教学内容的考核，难度和作业相当。考前可以复习一下上一周的作业。quiz 是闭卷，不能在桌上留东西，审核的不是很严，尽管如此不要作弊！如果有特殊情况没法考试要有假条，之后去助教那边补考。}
\SourceParagraph{期中 30\%，期末 40\%。期中期末基本上都是计算题，有一道证明题区分水平（一共十道题）。24 届考前会发考前练习，很有指导意义，基本考试的知识点、考点、题型就和考前练习差不多。考试题目以教材为导向。期末考试的难题证明题是书上的 problem plus 原题，看到就是赚到，考前还特意提醒我们去看 problem plus。考试可以带计算器，而且考卷上给了很多公式（如格林公式）。}
\SourceParagraph{最终这门课因为很多人成绩不太理想似乎老师还给全员抬了一下。}
\SourceLabel{微积分 II 学习建议：}
\SourceParagraph{微积分 II 大致对映国内的高数下，但不太可以用国内的教学平替，差别较大。这边仍然推荐 B 站 up 主“一高数”，但请选择性的看课，有的知识点不太能完全对应。}
\SourceParagraph{课本 james- Stewart- calculus 真的是很好的教材，请充分掌握课本上的题目。}
\SourceParagraph{微积分 II 的主要特点是套路性强，计算量大，知识方面是微积分 I 的加强版，且大多数定理证明复杂，所以反而不太强调证明，反而强调对于知识的运用。}
\SourceParagraph{十分建议以作业和 quiz 题目为导向出发学习，微积分 II 的重点在计算、应用、做题。微积分 II 可以看作是微积分 I 的延伸，微积分 I 没学好的同学要当心了。}
\SourceLabel{不挂科建议：}
\SourceParagraph{好好学。}
\SourceParagraph{作业是可以保证正确率的。挂科的同学大多数都是微积分 I 就没学好或者已经忘掉的同学。这边事先说明，如果没有微积分 I 的底子，微积分 II 学的是很艰难的。如果微积分 I 中的运算都会，微积分 II 大多数就是多重复几次你在微积分 I 中学到的运算，而且不用关注证明，只要会套公式和做模板题就行。请做到作业和 quiz 中所有计算题目都会算且有一个模板来计算这些套路题。没有自己总结能力的同学多用 AI 学习。虽然微积分 II 已经开始有难度，但是 AI 还是能处理的（但是 AI 可能会有计算错误，要用参考答案矫正）。考前练习的参考意义很大，如果实在来不及复习了就把考前练习喂给 AI，让他教会你每一道题目怎么写并且总结出套路，这样大抵能把考试里面的计算题都写对。}
\SourceLabel{高效卷绩点建议：}
\SourceParagraph{请把作业和 quiz 这两个确定且好拿的分拿到手，主要是通过对答案保证正确率。}
\SourceParagraph{以考前练习为导向发现自己有什么不足，考前练习的指导意义真的很大。作业和 quiz 虽然指导意义少一点但是知识点是相通的，也可以拿来复习。总之，对于微积分 II 这种强调计算，套路化的科目，以题目为导向学习是最快的。理论上只要你有题目和 AI 你练教材都不需要就可以考上 90 分了，把 AI 调教成你的私人老师。考试中有一些区分题很看微积分 I 的基本功，这里难以给出高效提高的指导，但是就算那些题目没做出来这门课考 90+应当也是没问题的。期末考试有教材上的原题，有高分追求的可以再花点时间把教材看一下。}
\SourceBlankLine[2]
\subsubsection{普通物理 I}
\SourceLabel{食用指南：}
\SourceLabel{1，这门课教学内容和老师短期内应该不会修改}
\SourceLabel{2，上半学期力学十分容易，但下半学期热学有点上压力了}
\SourceLabel{教学老师：}
\SourceParagraph{Sandro M. Ferreira Veiga（教力学，中文名张上梯）}
\SourceParagraph{Roope Halonen（教热力学，北欧大帅哥）}
\SourceLabel{教材：}
\SourceParagraph{力学有一本参考书叫 Matter and Interactions，但是书太厚了，老师也是按照 ppt 上课。}
\SourceParagraph{热学给了教材，叫 thermal physics，应该是赫大的教材，然后节选了一下}
\SourceLabel{课堂评价：}
\SourceParagraph{Roope 的课不听的话考自学很难精通热学，建议还是听一下的，不听课很难把握重点。而且 Roope 英文口语好，长得还养眼。力学是可以不听课自学的。}
\SourceLabel{给分：}
\SourceParagraph{力学和热学首先各占 50\%，每一门都是作业 30\%，考试 70\%}
\SourceParagraph{力学}
\SourceParagraph{作业 30\%。作业没有高考物理难，除了相对论那一部分正常人是理解不了的，建议直接背公式。如果你有物竞省二的水平都可以直接不听直接写作业。}
\SourceParagraph{考试 70\%。作业会写，ppt 全都看懂，考试至少 95+。考试和作业难度差不多，10 道大题，一体十分，24 届有一道推公式的题目，想拿高分的同学要能到做 ppt 中每一个式子都会推（相对论部分就算了）}
\SourceParagraph{热学}
\SourceParagraph{30\%作业。比较有难度，同学不会写，AI 也不会写。但是给分很好，全写完了，简单的题目不错基本都 95+。}
\SourceParagraph{考试 70\%。是可以带一面写满公式笔记的 A4 纸进考场的半考卷考试，比作业简单，有部分作业原题。带笔记的目的是为了让学生不要背公式。因为热力学只背公式而不知道怎么来和怎么用是做不出题目的，同时也是为了防止你名词不会拼或者英文不会拼。考试有一定创新性，所以很考验你对热力学的理解。考试有划重点，考前指导和复习课，基础题也很多，还有选择和填空，不用过于担心。}
\SourceParagraph{bonus point（附加分）：额外的分数，一共大概 7 分，还是很可观的。每周日前在邮箱中提交你预习课本的 reading diary，一次一分，态度分。}
\SourceListItem{24 届好像因为大盘数据不好看也提了一下分，大抵是 bonus point 发力了。}
\SourceLabel{普通物理 I 学习建议：}
\SourceParagraph{这门课因为国内没有很好的适配参考资料（如果非要的话可以看南京大学自己出版的《普通物理学》，但是那本书太厚了，如果你对自己比较有要求可以去看，体系和知识是想通的），所以没有任何网络资源推荐，这种时候 ai 就十分好用，不懂的概念，题目，知识点都可以问 ai。}
\SourceParagraph{力学篇建议有物理基础的同学真的可以直接看 ppt 推公式，ppt 都看懂了作业就没问题，作业没问题考试也就没问题。值得一提的是相对论部分，这部分除非你是天才否则人脑不可能完全理解，可以去先看看科普视频建立一个相对论时空观，然后背住公式，总结体型，会写题目就行，考试题目都是直接套公式的。}
\SourceParagraph{热学篇有点上难度了，难以速成且知识量大，知识点杂乱。这边建议首先跟上节奏，每节课都要好好听，这门课临时抱佛脚先抛开难度不说，时间花费是真的要死人的。这门课难就难在公式多，而且不会用，想要这门课考好就要掌握好公式，具体来说包括：知道这个公式怎么来的（具体推导不需要掌握，但是你最好大概知道是怎么推导的）知道这个公式什么时候用（就是知道这个公式是干嘛的，这一点做题多了就自然知道了），做到公式间的联系和区分（常常有一道题目要用好多公式，而且有很多计算同一个物理量但是根据不同的模型或不同的条件而适用的公式）}
\SourceParagraph{做题时依照这个逻辑；这道题已经知道什么物理量，要求什么物理量，要用什么公式（根据已知未知条件，或者做题多了以后的经验判断），计算（考试可以带计算器）。做到所有 ppt 都看懂，所有作业都会写，一定可以拿高分（但是这一点还是有点难度的）}
\SourceLabel{不挂科建议：}
\SourceParagraph{力学基本上就是高中的延申，如果你高中物理不差，力学不太可能挂，同时建议你不要力学尽量学高一点，因为热学可能会寄。热学这边建议你考前通过 AI 教学，教会你所有作业的写法，然后公式和名词尽量往 A 纸上抄，老师还有自己的总结，总之作业都会写的话应该不太可能会挂，把填空选择基本分拿到，作业也有原题出现。}
\SourceLabel{高效卷绩点建议：}
\SourceParagraph{力学自己看看 ppt，写写作业就能学的很好了。}
\SourceParagraph{热力学这边建议跟着老师的步骤好好来学，这边笔者本人虽然临时抱复脚成功了，但考试周复习时间成本很高。感觉让我再来一次会跟着老师好好学，而不是速成，但这边还是进行一些指导。}
\SourceParagraph{建议进行这样一个步骤速成：课本梳理知识体系（主要是知道学了什么）-做题（知道要考什么，怎么考）-把题目问 ai（因为这个时候你是不会写题目的，问 ai 这道题目的思路，解题步骤，对应知识点，然后让总结举一反三，争取做到这个题型完全搞定）}
\SourceListItem{-再去书上找到对应知识点得章节，看是否会遗漏。}
\SourceParagraph{老师自己有进行总结和考前指导，可以用那个来进行知识体系的建立和查缺补漏。}
\SourceParagraph{真的不建议考前速成，因为真的不高效，这门课考前速成至少 20h+}
\SourceParagraph{这里主要针对想拿96+的同学写一写考试中非常琐碎的地方；}
\SourceParagraph{期中是有回忆版的。老师考试之前会发复习题，群里有之前的复习题，请注意在本届复习题中被替换的题目。25届期中有一道往届复习题中的原题。}
\SourceParagraph{期末题型多样，又简答之类的，无论对题目多么一头雾水，请写点东西上去。两位老师批卷给分都好到令人匪夷所思的地步，就是会在你自己都想不到可以给分的地方硬塞分数给你。期末卷和期中区别还是很大的，期中和中学有点像，会推公式会算就行，期末最令我意想不到的地方是居然考了heat capacity,  thermal expansion coefffcient之类四个的名称，我感觉我期末扣分近一半都扣在各种名称上了（还有一个地方没想起来压缩是 compression），所以平时不要看到名称就跳直接看符号。另外最后几个应用Chapter本来以为不重要最后没想到考了一个大题，所以期末要考得好就不能抱侥幸心理。25届题干中出现humidty老师好像默认我们知道它的定义，但我记忆里没有学过这个东西。老师的靠前复习纯粹是把重点拎了一遍，没出现在里面的不代表不考，反复强调的也可能不考，和化学复习课完全是两个风格，没有一点透题的意思，所以全部都要好好复习。}
\SourceParagraph{另外“雨后竹林”里虽然是中文热力学的复习，但知识框架是一样的，里面讲的我们都学，对笔者个人把知识都串联起来十分具有帮助，而且没什么人听（我听的时候就三个人），对讨厌人群密集的人十分友好。}
\subsubsection{线性代数}
\SourceLabel{食用指南：}
\SourceParagraph{看课程体系是否更改，老师是否更改。如果更改较多，可能一些细致的总结就没啥用了。但是线形代数的自学资料是通用了，学习思路也是通用的（在做题中学习，以老师出题习惯为第一复习指南，善用 AI……）}
\SourceParagraph{教学老师：Daniel Mourad（不清楚是否会换老师，Daniel 是南大的博士生。）}
\SourceLabel{课堂评价：}
\SourceParagraph{讲课没有十分突出到让人觉得必须要听这个老师的课（）。且老师讲的内容在参考教材中都有。觉得浪费时间可以自学，想练听力可以听课。可以课间找 Daniel 问问题，人很好，不去找他问问题似乎还会略显寂寞。分数计算中没有课堂参与分，对社恐同学友好。}
\SourceLabel{给分：}
\SourceParagraph{课作/家作：35\%}
\SourceParagraph{课堂作业只有一两次，还可以带回去写，还可以重复提交。家作每次 20 分总分，给分基于正确性（一共 10 分）和完整性（一共 10 分，题目全都认真写了就给 10 分）。正确性主要是针对两三道难题和证明题做重点评价（不写过程直接出答案会在正确性上扣分，所以不能只写答案（教材参考答案很多情况只给出答案））。}
\SourceParagraph{作业用 canvas 平台提交（大部分设备要用梯子才能登陆）。canvas 平台似乎 ddl 错过了也可以提交，但是会有标注“迟交”。迟交似乎没事，只要在助教批改前提交，助教都是可以看到的。（我们这届有人晚了一个星期交作业仍然没事。但是要是你们自己要是因此惹出祸端，自己是第一负责人）}
\SourceParagraph{期中：30\%}
\SourceParagraph{期末：35\%}
\SourceParagraph{考试是不可以带计算器的闭卷考试。考试偏向基础概念的计算题，会有一道比较难的证明题压分，统计情况来看考试大于 90 分的人数在 10\%左右。作业都吃透考试就没问题。}
\SourceLabel{线性代数学习建议：}
\SourceParagraph{考虑到南赫不管怎么变改，一定都是基于外国教材进行学习，所以强烈推荐用国外（尤其是美国）的教学体系学习。书本建议使用 Gilbert Strang 教授的线形代数英语原版（MIT 的教材用书，同时也是 24 届的教材）。B 站和油管有这位教授的网课。附有配套的专门课程网站。}
\SourceParagraph{当然看不惯英文原版的也可以看中文译版，B 站也有用中文讲解这本书的，但是普遍播放量不高（例如 up 主“欧洲烤全羊”这样）。}
\SourceParagraph{同时可以先看一遍 3bule1brown 的详解线形代数动画，对于线形代数的几何理解会立即加深。}
\SourceParagraph{也推荐资料群中的 The-Art-of-Linear-Algebra，对于计算和矩阵分解会有更好的代数理解。}
\SourceParagraph{线形代数是一本工具和语言！可以从很多角度去理解（向量空间，方程组）。有要求的同学请不要拘泥于计算，请深入从不同角度理解。他是一门十分具有魅力的学科，且以后做科研可能要使用 matlap 就是专门处理线性代数的。}
\SourceParagraph{如果非要用国内的教材参考也可以，但是会面临教学顺序和考点不统一的困境（国内的教学强调计算和做题；外国的教学偏向理解和运用），从功利的角度来说浪费时间。但是在这里也做出推荐，可以看 B 站 up 主“当年线代”（可能对以后准备考研的同学有用？）}
\SourceLabel{不挂科建议：}
\SourceParagraph{上课可以完全不听}
\SourceParagraph{作业分请务必搞定（基本都可以 AI，如果是书本题请不要直接抄参考答案，参考答案很简略）}
\SourceParagraph{基础概念能搞清楚最好（可以通过上面推荐的非教材资料快速掌握，这个时候不建议用教材了，教材太厚了，还有不会的去问 AI），但是就算不是很清晰，也可以考前对症下药问 AI，做到作业中提到的名次词和算法会用。请做到作业中计算的题目全都能会写，不会的问 AI（60 分及格肯定是可以的）。}
\SourceLabel{高效卷绩点建议：}
\SourceParagraph{上课可以完全不听}
\SourceParagraph{作业请务必做到高分，把能确定拿到的分拿到手（可以 AI，可以找一个学霸同学）}
\SourceParagraph{考前利用好 AI，把考前练习每一道题目的知识点和方法都搞懂（重要！合理运用工具。找助教问问题不可能做到私人化直击痛点）}
\SourceParagraph{考前练习和作业在题型上和两次考试无强相关性。但是一点都没学的同学可以用考前练习作为复习起点（作业内容有点多，难以考前速通），通过做题和问 AI 题目高效学习，知道自己要掌握哪些知识点和做题方法（考前复习题可能是难于考试难度的）。}
\SourceParagraph{考试以计算题为主（60\%-70\%），少量证明题和一道有区分度的题目。要是作业或者考前练习里面每道题目的知识点都搞懂，且保证可以独立再次全部做出，考试 90 分没有问题（基础题基本可以全做对）}
\SourceParagraph{经过实践，24 届的副本可以仅依靠一周 1-2h 的作业时间和每次考前 2-3 天的复习坐牢速通拿到 90 分。}
\SourceParagraph{95+还是要做到课本全部掌握，}
\SourceParagraph{每周用时可能会膨胀到 5h 甚至 10h。}
\subsection{大二部分}
\SourceParagraph{这边建议多翻群里的资料，里面有上课ppt，教材和教材答案（很多时候就是作业答案）。并且群内会有学长学姐们背诵记忆版的考试题目，虽然每年要求有50\%的题目更换，但也同时意味着会有50\%左右的原题。就算没有原题也是高质量的复习材料，因为考察的知识和做题的方法是相通的。希望南赫大家庭的开源，共享，互助精神能一直传承。}
\SourceParagraph{PS:如果有些材料缺失可以问一下管理员/群主有没有备份}
\SourceBlankLine[1]
\SourceParagraph{大一大二的所有理科性质课程都由wq书写，本人所有理科课程均在90+，喜欢自学以及速通，且常常用ai学习（后面的ai学习部分也由本人书写）。撰写此部分是为了：1提供课程基本情况作为参考。2以拿到分数为目标，提供学习建议和指导。3提供以不挂科为目标的规划建议，提供高效率拿到高分为目标的规划建议。}
\SourceBlankLine[1]
\subsubsection{普通物理II}
\SourceLabel{食用指南：}
\SourceListItem{1. 这门课因为差评不断导致以后分为两个班教学，23，24届为徐海潭教授教学，且23，24届教学进度不一，仅作参考}
\SourceListItem{2. 在课程评价中打差评的都是学的不好的。笔者曾因为点评课程内容、说明老师教学良心被猛烈抨击。本人此门课总评96，一学期基本没有去听过课。并且这门课出分后好像平均分也在80分以上。可能后半学期的量子力学考的太难很多人破防了，具体情况后续会说。}
\SourceBlankLine[1]
\SourceLabel{教学老师：}
\SourceParagraph{徐海潭（很interesting，有自己的教学想法，但是可能对于一般同学来说会受不了，这门课他上的有种只为好学生开课的美感。有自己的公司）。}
\SourceBlankLine[1]
\SourceLabel{教材：}
\SourceParagraph{徐教授上课没有教材，没有ppt，根据自己的教案当场推导，当场讲。只能当场记笔记或者拍照。}
\SourceBlankLine[1]
\SourceLabel{课堂评价：}
\SourceParagraph{可以说这门课就不是为了一般人设计的，要你一学期学完人家物理系一年的内容,量很大，且没有任何学习便利（作业没官方答案，没有ppt）。老师这门课的设计有自己的想法。很看个人能力。与其纠结分数，不如把这门课当作一个测验，有能力和喜欢学这门课的人大概率是以后可以做科研的，因为这门课很考验智商或者耐心。笔者以及笔者的室友作为竞赛生都感觉真心不难。如果愿意花时间学，80分平均分还是能达到的。反正红不红，黑不黑大家都要学，不如来找学长学姐取取经。作业都会写考试就没问题，考试很多作业题。}
\SourceBlankLine[1]
\SourceLabel{给分：}
\SourceParagraph{具体算分不记得了，分为：作业+期中+期末+课堂互动。}
\SourceParagraph{作业是没有官方答案的，且有些问题是开放性问题，甚至有的问题会要解极其变态的常微分方程。请大家做好辨别，那些十分变态，看起来无法在考场上独立推导出来并且计算出来的题目是不可能考的，复习时不要抓错重点。}
\SourceParagraph{考试全是大题，每次考试都有20分附加题（就是说满分是120分，只要考到100分就算满分了），基础题基本都是作业题和上课例题。}
\SourceParagraph{课堂互动包含很多，回答问题和提问题都算，且是额外的分数（24届是每互动一次直接在总评上加1分，23届也有人没有互动拿到接近满分）。这门课老师是允许课堂上你不懂的内容直接问的，甚至可以用中文问（以后不确定是否要全英语）。要是自己听不懂不会可以直接和老师问一节课，因为老师自己只教一个班且自己出卷没有进度要求，教到哪里考到哪里（23届似乎出现了一直提问导致量子力学没有教多少的情况。而24届提问较少，量子力学都教完了。之后不清楚，因为变成两个班教学了）。}
\SourceParagraph{分奴的话就每节课疯狂提问，同时起到拖延进度+获得额外分（因为后半学期的量子力学很难）。群里面有学长记忆版本的考试考卷。}
\SourceBlankLine[2]
\SourceParagraph{电磁学，光学}
\SourceParagraph{前半学期基本上在讲电磁学和光学。本质上就是微积分的实际例子应用，然后模型和公式很多。电磁教学内容可以直接对应《新概念物理教程：电磁学》，编辑是赵凯华教授（群内有资料）。基本和第一章，第二章，第五章，第六章对应。光学似乎没有哪本教材完全对应，知识点比较分散，可以通过上网搜索，书籍拼凑，ai教学自行学习。}
\SourceBlankLine[1]
\SourceParagraph{量子力学}
\SourceParagraph{这门课真正重量级的地方。因为徐教授自己就是做量子计算的，所以说在这一块讲的十分变态，和匡院习题难度差不多。且量子力学这部分是根据徐教授自己的思路来讲解的，没有可以完全对照的教材，也不符合物理学史的发展来进行教学。但是考得都是平时上课讲的例题和作业里的内容，且教学的时候相同的内容会翻来覆去讲很多遍。}
\SourceParagraph{量子力学的难点在于：1数学工具很难，涉及到微积分，线性代数，泛函，模拟电路（量子计算）。2不符合物理直觉，不便于理解。3会有很多反直觉的案例和应用（量子隧穿，量子通信，量子缠绕，量子计算）。}
\SourceParagraph{这边建议一定要有一份完整的笔记。不一定要自己的，同学的也可以，作用是作为底板用来让ai教学。直接问ai对症下药，先学会其中的数学符号与数学工具，然后大致理解概念（概念理解这一步可以看一些B站上的硬核科普视频）。以上步骤是为了你可以读懂题目（理解概念）、会计算题目（数学工具），然后直接开始做题（作业于课堂例题），一开始完全不会可以让ai教学，最后做到能把题目全部独立做出的境界就可以，可以不用完全理解题目的背景和应用，会解题就行，期末绝对可以80+。虽然此方法不符合教授的开课初衷，但是绝对可以拿到平均分。（此方法推荐给上课跟不上或者想要速通的同学）}
\SourceBlankLine[1]
\SourceParagraph{最终这门课的平均分在80+，课程评价两极分化。徐教授是期望我们能成材，用高要求高期待培养我们的。教授私下人也很好、三观正，可以去约他聊天，受益匪浅。虽然有的时候会有抽象发言（例如说要把南大变成诺奖得主的摇篮之类的），但是对同学很好，所有问题都会一一解答。同学们不要有畏难情绪，如果能力足够，学好此门课，对于逻辑思维和数学物理能力是巨大的提升，同时对信心提升也很大（你都能自学量子力学了，还有什么学不了）；如果能力不足，即时和老师交流，善用机制，在投入时间的情况下也能拿到高分；如果彻底不感兴趣/觉得时间投入不值，成绩不是唯一的指标，学习理论知识也不是大学唯一的目的，摆正心态，不挂科，朝着自己的目标前进。}
\SourceBlankLine[1]
\SourceLabel{不挂科建议：}
\SourceParagraph{作业和课堂例题都会写的话应该不太可能会挂，在量子力学部分就算不能理解，也一定要做到能列出公式并且进行计算，平时分一定要拿满。}
\SourceBlankLine[1]
\SourceLabel{高效卷绩点建议：}
\SourceParagraph{使用ai学习，刚需有一份好的笔记记录下上课的所有知识点与例题作为底板。然后形成一个：提问——学习——实践（做题）——反馈——提问的学习循环。}
\SourceParagraph{这门课本人感觉不太能一两天速通。虽然笔者平日里没听课，但是私底下也是花了时间才学懂的，没有像别的学科一样能一两天速通（基本上量子力学部分花了一整周来潜心学习）。}
\SourceBlankLine[1]
\subsubsection{常微分方程}
\SourceLabel{食用指南：}
\SourceListItem{1. 此门课程大部分内容甚至考试题目基本取自于教材《常微分方程》，王高雄为编辑。}
\SourceListItem{2. 因为授课地点为苏州，且为专业内教授教学，短期内估计不会有更改。}
\SourceBlankLine[1]
\SourceLabel{教学老师：}
\SourceParagraph{王玉婷：上课基本上是在读ppt}
\SourceBlankLine[1]
\SourceLabel{教材：}
\SourceParagraph{有专门的ppt，但ppt内容基本照搬《常微分方程》（王高雄编辑，第四版）。甚至很多考试题目都取自于教材的例题和课后习题。群文件中有本教材和教材答案。}
\SourceBlankLine[1]
\SourceLabel{给分：}
\SourceParagraph{具体比例已经忘记，平时作业+期中考+期末考（好像可能会参照上课表现？）。}
\SourceBlankLine[1]
\SourceLabel{课堂评价：}
\SourceParagraph{王老师上课基本都在读ppt，会有一些小互动拉同学上去解题。但是ODE这门课本来就是一门应用和应试的课程，也讲不出什么花样。这门课的学习就是掌握方法，然后刷题，其中涉及到一点点的背公式和推导（主要集中在迭代逼近那边），考试的题目基本上全是解方程计算题，然后来一些神秘情景题（24届有一道题的背景涉及到了化学一级反应和二级反应的概念，导致很多人不会写）。}
\SourceBlankLine[1]
\SourceParagraph{感觉是大学学的理科中最不需要脑子的一门课，但是对于熟练度要求很高，否则考试题目做不完。感觉这门课只有不及格/平均分/高分三种状态，不及格的就是基本上解题方法没掌握或者分不清；平均分就是把作业都会写了，但是熟练度不够且难题做不出来；高分就是把《常微分方程》（王高雄编辑，第四版）全部做完，然后研究其中的一些偏难怪题（基本上每次考试会有一两道非常规的题目，但是基本上还是在书中的习题范围中，毕竟ODE的题目就只有那么多，暴力一点全部做完就行。但是没有完整做出来不影响拿90+）。然后这门课对于微积分的要求很高，建议复习一下微积分，特别是一些三角函数/反三角函数怎么积分，分部积分法怎么用。}
\SourceBlankLine[1]
\SourceParagraph{建议这门课多刷题，把《常微分方程》（王高雄编辑，第四版）上面所有题目刷掉就行，如果还不放心可以在群文件里找到以前的习题课ppt，里面有一些比较变态的题目可以尝试一下。}
\SourceParagraph{速通也是刷题，量只要把课堂ppt里面的题目和作业题都会做就行了。}
\SourceParagraph{但是就算把题目全部做完好像也就一个星期的时间。}
\SourceBlankLine[1]
\subsubsection{概率论与数理统计}
\SourceLabel{食用指南：}
\SourceListItem{1. 此门课前后半学期老师不一样，且统计和概率是两门学起来不太一样的学科。}
\SourceListItem{2. 这门课可以申请不参加期中，只参加期末，然后期中的分数比例全部用期末成绩计算}
\SourceBlankLine[1]
\SourceLabel{教学老师：}
\SourceParagraph{王伟：教学概率论，有ppt}
\SourceParagraph{蒋新宇：教学统计，课程ppt在自己的专门网站上，里面有一些互动小游戏，此老师还有自己的出版物与公众号（“人间的统计学”，公众号为“大家的风信研”）。上课有的比较难的部分会用中文讲。}
\SourceBlankLine[1]
\SourceParagraph{教材：《Probability and Statistics》Fourth Edition（作者：Morris H. DeGroot}
\SourceParagraph{CMU Mark J. Schervish CMU）教材主要是提供作业题目。群文件中有教材和教材答案}
\SourceBlankLine[1]
\SourceLabel{给分：}
\SourceParagraph{作业+期中+期末：20\%+30\%+50\%。可申请不参加期中，只参加期末，然后期中的分数比例全部用期末成绩计算。期中期末都是只有大题}
\SourceParagraph{期中只考概率论，大部分是计算题/应用题，会有少量证明题（难度高的题）。}
\SourceParagraph{期末会考察一半的概率论，一半的数理统计。其中概率论偏向计算题，数理统计偏向应用题（数理统计本来就是强应用性的）和一些神秘的名词解释，流程解释（考察对于数理统计流程的理解和关键定理的理解）}
\SourceBlankLine[1]
\SourceLabel{概率论：}
\SourceParagraph{概率部分偏向理论数学，需要一定的组合数学知识和微积分知识。对于一些概念的理解和定理的使用条件一定要理解到位，且期中考试中有一定量的证明题。期中内容全部是概率论，且普遍反应比期末难。但是证明题做不出来并不影响你总评拿90+，因为期末会比较简单。如果想要做出证明题建议一定要把ppt和书本内容吃透。但是作业题都会写的话期中拿到80+应该不是问题。}
\SourceBlankLine[1]
\SourceLabel{数理统计：}
\SourceParagraph{数理统计偏向应用，不一定要完全掌握推导过程和证明，但是对于模型的辨识，定理的理解应用，公式的背诵要十分熟练（特别是公式的背诵，在后面的假设性检验，卡方检验都要背住）。在期末的考察中只要能把ppt中的例题全都会写，相关的应用题就绝对没有问题。总的来说比前半学期好学（就是对于公式有一定的记忆要求，或者说你能全部推导，笔者这边就没有死记硬背，而是通过基本公式来进行推导并且记忆）}
\SourceBlankLine[1]
\SourceParagraph{强烈推荐自学和速通的同学看B站up“一高数”的概率统计速通课}
\begin{figure}[H]
  \centering
  \includegraphics[width=0.92\textwidth,height=0.78\textheight,keepaspectratio]{figures/wq\_image1.png}
\end{figure}
\SourceParagraph{之前微积分和线性代数好像也推荐过这个up的课程。只要是大学的数学，这位博主的课程简直是神。只要能花足够的时间把其中的内容吃透，不挂科甚至考到满分都是可能的。但是注意这位up的课程未能对考点全覆盖（我们概率论内容最后一张在这个B站课程里是没有讲的）。但是只要囊括了80\%的考点。只要能把这个看会是绝对不挂科的。如果再搭配上课堂中的ppt和作业，概率论与数理统计是十分好学的。}
\SourceBlankLine[1]
\SourceParagraph{只要能把这个课程看完。我保证你绝对不会挂科。}
\SourceBlankLine[1]
\SourceParagraph{高效卷绩点的同学建议把速通课看完，然后把ppt里面的概念做到全理解，题目全会做，就没有问题了。}
\newpage
\section{方向介绍}
\SourceParagraph{由于南赫学院专业不是很全，所以这里从大气科学学科总体来介绍。}
\SourceParagraph{大体来讲，南大大气主要分为三个部分，气候，中尺度，大物大化。这种划分是非常不具体的，每一个部分都可以进一步分为几个子部分。目前来说，几乎没有人能掌握大气的所有领域，所以笔者也只是根据个人了解提供一部分方向的介绍。如果有未能覆盖到的方向，还请各位读者自行查阅老师的履历，通过邮件、面谈等方式获得进一步的认识。}
\subsection{大气动力学}
\SourceParagraph{大气动力学是最近一百年里推动大气科学产生质的飞跃的方向。这一名称作为研究方向是具有迷惑性的，因为研究大气动力学的学者的实际研究内容可能是多种多样的，这并不是一个具体的研究方向。所谓“动力学”，指的是研究运动发生的机理的学科，也就是“力”的机制的学科。在大气领域，这种“运动”可以是大尺度的波动、涡、环流，可以是中尺度的气旋、对流、锋面，也可以是更小尺度的局地地形波、边界层等。“动力学”的目的就是揭示这些运动背后的物理过程。}
\SourceParagraph{上世纪三四十年代开始，以 Rossby 为代表的一系列科学家开启了大气物理学的近代化，即大量引入数学、物理方法解释大气现象。这就是芝加哥学派。芝加哥学派的特色是在准地转的框架下考察各种波动和不稳定。通过对不同的研究对象进行不同的化简，他们考虑了正压斜压、经向纬向、有粘无粘等条件的组合，然后在此基础上对基本方程组（NS 方程，质量守恒，能量守恒等）进行化简，能求解的线性波（不稳定）直接求解，不能求解的弱非线性波则引入摄动法（线性化），成功解释了大量中大尺度的天气现象。到八十年代，芝加哥学派集大成之作，Pedlosky 的”Geophysical fluid dynamics”，以及动力气象的经典教材，Holton 的”Introduction to dynamic Meteorology”也是这一时期的作品。我国大气科学界深受芝加哥学派影响，不仅大量引入相关文献书籍，更深度参与其中，其中不乏知名者，如叶笃正，郭晓岚，顾震潮，曾庆存等，出版书籍如《大气环流的若干基本问题》《数值天气预报的数学物理基础》等。这是一个出大师的时代，四海翻腾云水怒，五洲震荡风雷激，这一套理论配合上大师们娴熟的数学方法，迅速求解了大量微分方程组，解释了大量大气现象，奠定了现代气象学的基础。然而，芝加哥学派的研究方法在大师将其整理成书那一日起，就已经意味着自身的式微，半个世纪的不懈求索中，其已经解决了力所能及的容易解决的几乎所有问题，新研究者逐渐减少，在新世纪，风头更是被全球变化等问题抢去。}
\SourceParagraph{与此时间稍晚的是英国的剑桥学派，比起求解方程，他们更关心那些方程的几何属性，使用几何来理解天气系统。剑桥学派对数学工具并不挑食，比较关注强非线性的流体力学难题，相比之下，进入新世纪以来仍然比较活跃。}
\SourceParagraph{当然，上面说的只是动力学的冰山一角，根据研究对象划分，还可以有气候动力学，台风动力学，对流动力学，边界层动力学，地形动力学，热带动力学等等。从学科交叉的角度来说，比较相似的还有海洋动力学，行星动力学等，可以考虑了解。}
\SourceParagraph{气候动力学，研究对象为时间和长度尺度比较大的对象，比如大气环流、急流、涡流相互作用等。当然研究气候系统的也是有的，研究对象如 ENSO 等。从动力气象的角度看，更相近的是前者，南大在这方面代表性的老师是大气院的张洋老师；后者的代表性老师是南赫的孙德征老师（孙老师目前的重心在搓气候模型）。}
\SourceParagraph{台风动力学，主要针对台风各生命阶段的机理，如生成、增强、耗散等，谈哲敏校长大组是南大研究台风的中心，除了他本人外，唐晓东、仇欣、李元龙等老师也从事台风动力学研究，南赫方面，陈婕老师研究台风登陆过程中的动力机制。据李元龙老师所说，虽然使用数值模式我们已经能够很容易模拟出台风，但是我们并没有比较好的理论从基本方程组出发解释台风，因此许多研究其实还是在提出假说的阶段。}
\SourceParagraph{对流动力学，研究的是各种类型的对流，如浅对流、深对流、组织对流（超级单体、飑线等）等，以及相关的对象如积云等。这方面研究的流体力学成分多一些，因此可以使用实验模拟或者类比。对流动力学和台风动力学是非常接近的研究方向，因此有些老师两方面都会做，如顾剑锋老师、杨博雷老师，研究对流动力学的老师还有傅豪，他同时还做流体实验。}
\SourceParagraph{边界层动力学，包含 Ekman 层动力学，是与湍流息息相关的研究方向。现在活跃在大气边界层动力学的老师有周博闻等。}
\SourceParagraph{地形动力学，包括山地气象，与边界层动力学息息相关，都是中小尺度的动力学，因此二者重叠颇多。这一方面纯理论工作能做的不多，更多的还是在原有理论的基础上结合现实观测做一些解释和针对性的改进。这方面主要是徐昕老师在做相关工作。}
\SourceParagraph{地球流体力学（GFD），是包含大气动力学的更广阔的研究范畴，包括海洋动力学（研究 Overturning/洋流/大尺度漩涡等）和行星科学（其他行星大气结构，如急流、涡旋、超声流等，也有冰卫星等）等，都是同一套方法论（旋转系中的 NS 方程等），目前算是新兴、前沿的方向。许多理论是同时运用于不同的对象的。可以说，是在广阔的参数空间中进行研究，而不局限于地球大气。对更一般的物理规律而不局限于气象或者大气的同学可以考虑一下这个新兴方向。目前算是在做 GFD 的老师，广义上讲，做得比较理论的都可以算，如果狭义上说，则是以傅豪为代表的老师。需要注意的是，普林斯顿和 NOAA 有个 GFD 实验室（GFDL），南赫也有一些老师在那呆过（比如孙德征、陈婕等），但是他们的工作并不是很纯的 GFD，而是什么方向都可能有的。南大跨院系合作比较少，如果对这一方向感兴趣也可以找其他院系的老师，比如目前行星科学里做 GFD 的有深空探测研究院（属前沿科学学院）的倪冬冬老师，海洋动力学则比较少，笔者尚未听说有做这个的（欢迎补充）。}
\SourceParagraph{总体来说，目前专门从事动力学研究的学者还比较少，并且不是热门方向。上世纪动力学的大发展让这门学科的理论深度大大提高，在有些方向上停滞不前，在有些方向上又太过深入，在有些方向上又太过无力，于是要么前途未卜，要么门槛高，要么无从下手。因此，只建议有兴趣，有理想，数理基础扎实的同学专门从事这方面工作。}
\SourceParagraph{需要注意的是，虽然大气动力学听起来做的工作是物理学家或者数学家的工作，但其实不然，由于无法真的做实验，因此从事大气动力学工作的学者往往也会从事以下工作：数值模拟（必备技能，不会搓模型也得会跑模型/调参，能看懂别人的模型），数据分析（观测结果），实验（类比）。纯粹理论研究也很难发文章，因此理论研究者也需要参与到应用工作，比如参数化，比如服务国家需要的各种项目等等。}
\SourceParagraph{说了这么多，如何入门大气动力学呢？通常根据具体的方向有不同的学习路径，但是在本科阶段共同的路径还是有的：微积分+线性代数+常微分方程+流体力学+动力气象+地球流体力学。这里重点推荐一下动力气象和地球流体力学。动力气象读 Holton 那本 Introduction to dynamic meteorology，就能大致了解上个世纪九十年代之前的主要进展。但是这本书中文翻译一大坨，英文写得其实也没那么好读，相较之下也没有什么更好的教材，如果未来出现了的话可以再推荐推荐。地球流体力学比较推荐 Vallis 写的教材 Atmospheric and oceanic fluid dynamics，这部教材相对而言现代一些，是 Vallis 根据他在 90 年代到新世纪初这段时间的教学实践总结的，实际上融合了大气动力学、大气环流和海洋动力学，经常被作为大气环流/地球流体力学的参考书。这书很厚，可以按需挑着阅读。再往后就看方向了，对波流相互作用，湍流，边界层，各取所需，有些还有教材，更进一步的就只能看综述性文献了——事实上，能到啃下当代综述性文献的水平，基本就到了从事动力学研究的门槛了，但是光达到这一点就需要耗费大家大量的时间和经历，所以我用黄埔军校的对联来建议大家：“升官发财请往他处”。决心进入这一领域的同学，也要学习前辈的治学精神，赵九章，叶笃正，曾庆存，伍荣生，乃至我们的师长，如谈哲敏、王元、周博闻、顾剑锋、傅豪等人。南大的动力学基础已经是国内顶尖的了，大家要利用好两院的资源。另外，对于想从事这一领域工作的同学，建议数学物理可以学深入一些。南赫学院没提供的课程，或网课，或旁听，或读书，多少得接触一些。数分高代，复变实变泛函，数理方法，理论力学，热统，甚至电动力学最好也接触一些，对研究流体也是有帮助的。大气动力学作为物理学，或者说流体力学的分支，不应该对数学的、物理的方法挑食，时间太紧的话，就不必要修课了，能啃下来就好，不要把时间花在刷题和准备考试上，细枝末节的问题也不必死磕。另外，黄思训教授和伍荣生院士曾写过一本《大气科学中的数学物理问题》，大家学完一些数学物理课程之后，可以看看这本书，结合大气中的问题重新审视一下那些数学物理技巧的用途。}
\SourceBlankLine[2]
\subsection{大气物理}
\SourceParagraph{由于南赫目前还没有很多对大气物理（云物理/辐射等）方向有足够认识的本科生，故采取约稿的方式为打算走进大气物理领域的同学们提供一些这个方向的学习体验。撰稿人是大气院 18 级的风云人物，这篇短文的文字和组织也比较有趣，相信大家从中能得到一些对大气物理方向的思考。}
\SourceParagraph{撰稿人：msk（18 级大气本，22 级大气 PhD）}
\SourceParagraph{「大气物理和大气化学是大气科学大厦上的两朵乌云。」}
\SourceParagraph{大气物理是一个筐，什么都能往里装。翻开那本经典的「Atmospheric Sciences: An Introductory Survey」，除去其中五个分别关于大气热力学、大气化学、大气动力学、天气系统与气候动力学的章节，剩下五个章节全都是大气物理学的内容。但是，大气物理学的问题看起来不像传统大气科学那么“理论可解”，例如：京都的城市建筑如何影响大气的流动？挪威的森林如何为大气增加湿度？加州山火排放的气溶胶如何影响云的形成、又怎么影响大气的辐射特性？等等。}
\SourceParagraph{这是由于大气物理学与传统大气动力学、气候动力学最显著的差别，在于大气物理学没有一个兼顾可靠和实用的“第一性原理”。在大气动力学中， Navier-Stokes 方程是当之无愧的金科玉律——如果你没能解释一个动力学过程，那一定是你没有正确地对 N-S 方程进行变换、拆分和重组。但大气物理学研究的就是 Navier-Stokes 方程里面那些令人讨厌的「非绝热加热项」和「湍流混合项」。我怎么用方程写出：屋顶上每个瓦片对湍流的影响？树上每个叶片光合和蒸腾的速率？每个气溶胶粒子可供凝结云滴和核化冰晶的比例及其辐射属性？更何况，城市由无数瓦片构成，森林由无数树叶构成，野火的烟雾里气溶胶粒子的数量更是那由他！}
\SourceParagraph{在学习大气物理的过程中，需要时刻告诫自己：人类其实知道得很少。很多模式中的物理方案都只是凑合着能用，并且缺乏从第一性原理改进的空间。在基于模式的天气预报和气候预测中，动力学上的误差可以归咎于网格不够密，但大气物理的误差只能归咎于人类还没有发展到那一步了（笑）。令人沮丧的是，大气物理方向的毕业生通常比天气、气候方向的毕业生更难在国内的气象系统找到工作，或许是因为大气物理离实际的天气预报太远、又或者是大气物理学子对自己的工作是否真正能改善天气预报持怀疑态度吧。}
\SourceParagraph{但同时，这是一个时刻有新知识被生产出来的方向。人类逐渐认识到自身的发展与气候变化、地球系统之间有千丝万缕的关系。大气物理、大气化学目前还处于一个百花齐放、蛮荒发展的状态，新的文章数量开始井喷。虽然大气物理方向的研究还是与动力学脱不开干系的，但大气物理这个大筐内部隔行如隔山。对于新踏入大气科学殿堂的大家，我的建议还是按照自己兴趣来选自己研究的方向——毕竟对着一个谁也不知道靠不靠谱的东西，如果自己没兴趣还是比较痛苦的。对于笔者来说，那本经典的「Atmospheric Sciences: An Introductory Survey」对每个大气物理方向都是一本很好的导论，还是推荐大家对感兴趣的章节仔细阅读。笔者文笔不好，那就先这样（}
\subsection{大气电学}
\SourceParagraph{南大本没有做大气电学的，但是浦云娇老师加盟南赫学院之后，我们也算是有了这一方向。下面的介绍由傅豪学长撰写，浦云娇学姐校对，大家注意一下时效性。}
\SourceParagraph{大气中有丰富的电磁现象：与积雨云相伴的闪电、高层大气的极光，以及一种从对流层贯穿到电离层的瞬态放电：红色精灵 (red sprite)。}
\SourceParagraph{我们所处的对流层是电中性的，所以电磁过程对天气、气候的长期影响主要通过影响大气化学来实现。比如，闪电可以促进氮气转化为氮氧化物。闪电对古气候也有特殊意义。在约 40 亿年前，原始大气(proto atmosphere)形成的过程中，大气的快速冷却导致大量水汽凝结，海洋被认为形成于这场持续数百年的暴雨。该过程中频繁的闪电被认为对生物大分子的形成有重要贡献[1]。}
\SourceParagraph{大气电学最迷人之处还是其物理本身。大尺度高层大气运动会受电磁场影响。在电离层(ionosphere)，稀薄大气被太阳辐射加热，在白天部分电离、夜晚恢复中性。电离层大气运动受低层大气上传的重力波、Rossby 波和太阳辐射加热共同驱动，自身除了受旋转、层结影响外，还受到洛伦兹力。电离层大气动力学尚不成熟。其大尺度运动被认为可以用“磁浅水方程组”(MHD shallow water equation)来刻画[2]。}
\SourceParagraph{在 小 尺 度 上 ， 大 气 电 学 和 云 物 理 有 许 多 交 叉 。 一 个 问 题 是 云 起 电 (cloud electrification)，就是冰晶碰撞如何导致正负电荷分离，可以在云室中(cloud chamber)研究[3]。另一个是气体放电(gas discharge)，涉及等离子体物理学(plasma physics)和高能物理(high-energy physics)。人工触发闪电是个很有趣的问题，可以帮助人们更可控地观察这种瞬间发生的过程。比如火箭触发[4]、激光触发[5]。}
\SourceParagraph{大气电学研究方法包括观测、数值模拟和少量室内实验。对雷电的观测，可以通过高速摄像机捕捉光学信号，或者用天线阵进行辐射源反演。观测需要很多动手操作、甚至仪器设计，可慢慢培养出电子工程素养。对气体放电的数值模拟没有像 WRF 那样的成熟的软件，有时要自己动手写计算模型。室内实验主要和高电压工程交叉，比如清华电机系的气体放电实验室。}
\SourceParagraph{总而言之，大气电学在大气科学中很小众，但和等离子体物理、电子工程有许多交叉，可能一不小心就离开大气了。对物理有执念的同学可以考虑一下。以下是笔者推荐的课题组：}
\SourceListItem{（1）郄秀书研究员,中科院大气物理研究所。从事人工引雷、闪电发展物理过程的研究，偏观测。在山东沾化建有引雷基地，在那里观测很有趣。央视拍过专题片介绍：}
\SourceParagraph{\url{http://tv.cctv.com/2018/01/09/VIDEBtDZkhvl0tBErseSNOOJ180109.shtml}}
\SourceListItem{（2）Steve Cummer, Duke University, 电子与计算机工程系。电磁波专家，关注如何用天线阵反演闪电精细结构、闪电的高能辐射现象。主要做观测、仪器:}
\SourceParagraph{\url{http://people.ee.duke.edu/~cummer/}}
\SourceListItem{（3）Victor Pasko, Penn State University,电子工程系。优秀的理论家。有水滴放电、电 离 层 物 理 等 多 个 方 向 。 研 究 手 段 为 数 值 模 拟 :}
\SourceParagraph{\url{https://www.eecs.psu.edu/departments/directory-detail-g.aspx?q=VPP1}}
\SourceListItem{（4）Vladimir Rakov, University of Florida，电子工程系。全方位的闪电专家，偏重雷电防护、人工引雷: \url{https://www.ece.ufl.edu/people/faculty/vladimir-rakov/}}
\SourceListItem{（5）Ningyu Liu, University of New Hampshire，物理系。优秀的理论家。对于成绩不那么突出，但真心想搞基础理论研究的同学，可能是个特殊的机会。}
\SourceParagraph{\url{https://mypages.unh.edu/ningyuliu}}
\SourceListItem{（6）J. R. Dwyer, University of New Hampshire，物理系。高能物理理论家。和 Ningyu Liu 经常一起合作。感觉他俩课题组是连通的。}
\SourceParagraph{\url{https://findscholars.unh.edu/display/jrd1002}}
\SourceLabel{另有两位国内老师以大气电学为副业，供参考：}
\SourceListItem{（1）武慧春教授，浙江大学物理系。从事等离子体物理研究，大气电学为副业。一不小 心 对 球 状 闪 电 产 生 兴 趣 ， 提 出 了 一 个 有 趣 的 理 论 模 型 :}
\SourceParagraph{\url{http://ifts.zju.edu.cn/content/?341.html}}
\SourceListItem{（2）庄池杰教授，清华大学电机系。从事气体放电的室内实验研究，高电压工程专家，不直接研究闪电，但做的方向在物理本质上很接近:}
\SourceParagraph{\url{http://me.tsinghua.owvlab.net/virexp/sys/gdys/}}
\SourceLabel{参考资料：}
\SourceListItem{[1] Chameides, W. L., and James CG Walker."Rates of fixation by lightning of carbon and nitrogen in possible primitive atmospheres." Origins of life 11.4 (1981):291-302.}
\SourceListItem{[2] 磁 浅 水 方程:\url{https://www2.hao.ucar.edu/sites/default/files/users/sheryls/cho-gtp16.pdf}}
\SourceListItem{[3] Ávila,Eldo E., et al. "Charge separation in low-temperature ice cloud regions."Journal of Geophysical Research: Atmospheres 116.D14 (2011).}
\SourceListItem{[4] 郄秀书,et al. "新型人工引雷专用火箭及其首次引雷实验结果." 大气科学 34.5(2010):937-946.}
\SourceListItem{[5] "激光引雷研究中的若干基础物理问题." 高电压技术 34.10 (2008).}
\subsection{行星大气}
\SourceParagraph{众所周知，大气科学既可以按照时空尺度分为微观（如大物大化）、小尺度（如边界层）、中尺度（如组织化对流）、大尺度或行星尺度（如行星波、气候），也可以按照研究对象或研究方法分为大气动力学、大气物理、大气化学等。上述方向或多或少会与其他学科交叉：既有与环境科学、化学交叉的，也有与物理学、数学交叉的，还有与数据科学、计算机科学和人工智能交叉的。然而，在传统研究视域中，与天文学交叉的研究较少。换句话说，即使我们整合了上述方向，研究的是地球系统科学，也没有跳出“地球”这一范畴。而行星科学则将研究目标拓展到其他行星，行星大气研究自然关注其他天体上的大气。这一领域乍一看应该完全属于天文学的研究范畴，但研究其他天体的大气仍需以已有的地球大气研究成果为基础，因此，可以将行星大气研究视为把地球大气科学拓展到一个极为广阔的参数空间。过去我们所做的各种假设和约束都可以被打破，或许是化学成分，或许是行星尺度的参数变化，甚至需要考虑行星轨道参数对大气的作用。这一点与古气候研究颇为相似。这样一来，我们对外部世界了解得越多，同时也就意味着我们对地球大气的理解也会越深刻，因此，行星大气和地球大气研究是相互促进的。下面，我将先介绍行星大气的研究概况（注意，我所提及的科学问题和科学发现并不全面，仅供抛砖引玉，实际上可供研究的课题还有很多），然后再介绍一下我对目前国内外行星大气研究的了解。}
\SourceParagraph{行星大气研究和地球大气研究的最大区别在于原位探测极为困难，所以在火箭能够将探测器发到其他行星轨道之前，我们对其他行星大气的研究基本上依赖望远镜观测，而望远镜的放大倍率又是有限的，因此我们最多能够看到行星尺度的情况，比如木星上的大红斑和急流。六十年代之后，卫星上天，美苏航天争霸让行星科学家有机会近距离研究行星大气，具体来说，我们可以通过飞掠、环绕、着陆等方式获取行星信息，比如大气成分、大气环流、温度廓线等。对太阳系内行星的探测在随后半个世纪中迅速展开，取得了大量进展：}
\SourceListItem{· 金星}
\SourceParagraph{金星大气对可见光几乎是不透明的，所以在望远镜下几乎看不到具体细节，但是上世纪有大量的探测器抵达金星并着陆，这让我们对其温度廓线、成分分布都有了大量的认识。然而这种研究主要集中在上个世纪，上个世纪甚至曾有两次在金星上释放气球进行探测；近二十年间，在轨探测任务主要有 ESA 的金星快车（Venus Express）和 JAXA 的晓（Akatsuki），但目前均已退役，因此近几年没有金星探测器在轨运行。不过，进入 2030 年代后，相关任务有望逐渐增加，NASA 正在计划的任务有达芬奇（DAVINCI，计划 2030 年发射，2033 年进入金星）和真理（VERITAS，计划 2031 年发射），ESA 正在计划的有 EnVision（计划 2031 年发射），我国的金星探测任务计划（2033 年前后发射，计划为金星大气潜在微生物采样返回，目前中央尚未批复），据说印度和俄国也有计划。}
\SourceParagraph{金星地表气压为 90 bar，温度达到了数百 K，大气成分主要为二氧化碳，其次为二氧化硫和氮气等，也有少量水，对流层延伸至 100 km，云层约在 50 km 以上，但是以下为雾霾，所以地表能见度都不高。从位温廓线看，地表数千米有显著不稳定，一般认为是成分梯度导致，此为表面混合层，往上基本都是稳定或中性的，其中在 20—30 km 上下存在一个中性层（deep-atmosphere mixed layer），50 km 左右存在中性层，这一层的对流很容易触发，被称为 middle-cloud convective region，不过对流基本都是浅对流，主要也就起到搅拌大气的作用。}
\SourceParagraph{金星大气最主要的科学问题有两个，一个是超旋转（Superrotation，也就是说仅用角动量输送无法解释的强急流），一个是大气物理和化学过程，包括探究其温室效应、潜在微生物等。目前这些问题都有文章探讨。}
\SourceListItem{· 火星}
\SourceParagraph{火星探测器的数量比金星探测器多不少，使用寿命通常也更长，我国的任务有天问一号（祝融号火星车）和天问三号（计划 2028 年前后发射，中央已批复，主要面向取样返回），美国也有好几台火星车，这些更加知名，此处不再详述。}
\SourceParagraph{火星地表气压只有 6 hPa，也就是 0.006 bar，99\%为二氧化碳，其余为氮气等，也有少量水。火星大气环流的特点是，相比地球，Hadley 环流只有半边，因为南半球整体比北半球高，所以上升支总是在南边，下沉支总是在北边，而非近似赤道对称的环流；与地球相比，季节对火星大气环流的主要影响是增强或削弱环流，而不会造成环流的跨半球移动。火星上最显著的天气事件是全球性沙尘暴，几乎隔几年就要来一回，沙尘几乎覆盖全球，除此之外，在小尺度上还有 dust devil 这种类似尘卷风的现象也频繁出现，目前认为与低层对流有关。}
\SourceParagraph{火星大气最主要的科学问题就是全球性沙尘暴的预报，现在国内有好几家机构在做火星大气模式，比如中科院大气物理研究所、澳门科技大学等。火星地表还有二氧化碳冰和水冰，以及二氧化碳云和水云，其分布、形成、迁移也有大量研究。需要注意的是，虽然火星上有水冰，但是水冰被太阳辐射加热后会直接升华，因此没有液态水。}
\SourceListItem{· 木星与土星}
\SourceParagraph{对外太阳系的探索要少很多， 最先对气态巨行星进行飞掠探测的是先驱者 10 号和 11 号（Pioneer 10 \& 11），旅行者 1 号和 2 号（Voyager 1 \& 2），留下的资料主要是照片、低阶引力场参数、 温度等， 以及其卫星的近距离观测。 后来木星有了伽利略号 （Galileo） ，向木星投入了进入器研究其垂直结构，以及朱诺号（Juno），土星有了卡西尼-惠更斯号（Cassini-Huygens），这三个探测器对木星和土星进行了前所未有的深入研究，极大加深了我们对木星和土星大气、内部结构、卫星的认知。目前，ESA 发射了冰卫星探测器 （JUICE， 2023年发射， 预计2031抵达） ， NASA发射了欧罗巴快船 （Europa Clipper，2024 年发射，计划 2030 年抵达），这两者的主要目标是冰卫星，重点包括木卫二和木卫三，我国也在计划天问四号（2029 年前后发射，中央已批复，主要环绕木星以及研究木卫四）。}
\SourceParagraph{木星和土星的大气都非常深厚，在光学上是不透明的，因此我们对其内部知之甚少，引力场和电磁场的探测结果也基本集中在上层，但是仅看其外层就能够获得大量至今未研究明白的大气现象，比如大红斑（巨大的反气旋，最近几十年的观测显示正在持续缩小，目前认为其起源是行星尺度的不稳定遗留下的涡旋）、急流（这两颗行星也有超旋转现象，以及最重要的是流向交替的条带状急流）、极涡（对气态行星极涡的系统研究也就是最近十年的事情，因为之前没有探测器完整看到过木星极地！）等等。气态行星上的化学反应、天气层与内部的耦合等也都是目前研究的热点。}
\SourceParagraph{需要注意的是，气态行星的内部热通量比接收的太阳辐射通量要多，但是如何将这一事实完美地与模式结合，目前还没有定论。理论上可以在模式中利用完整的辐射模型，只是辐射源的方向反过来，但是相关研究尚未充分开展（主要是缺人缺钱）。如果要总结木星的特点，可以概括为以下三个方面：第一，木星是层结稳定的，但是中高纬度有大量的强烈对流（即可以观测到闪电，这里有个很有意思的点是木星低纬的闪电非常少，中高纬比较多，因此普遍认为木星极地的热通量要高于赤道）；第二，木星上的示踪物有垂直梯度但是没有水平梯度；第三，木星上的熵在水平和垂直方向都有梯度。}
\SourceParagraph{目前研究气态行星大气环流所用的 GCM 主要分两派，一派源于地球 GCM，一派源于恒星大气模型，二者各有优劣，目前认为结合二者才能比较好地模拟气态巨行星大气环流，因此做行星大气最好能够对恒星大气有些基本的了解，互相借鉴。由于卡西尼号已经坠入土星了，目前气态巨行星上的环绕器只有朱诺号。朱诺号目前的轨道主要服务于极涡研究，李成老师引领的相关团队正在利用逐渐升高的朱诺号轨道对木星极地进行更详细的研究，尤其是利用微波波段成像研究其垂直结构。目前研究气态行星内部的手段非常有限，除了探测器，最经济也最直接的方式是有个彗星砸进木星，从而可以测量其不同层的 Brunt–Väisälä频率，以及内部成分等，当然这种事情难以预期，上次发生在 1994 年；近几年有一个非常有意思的方法，就是利用行星流体内部波动在行星环上引发的涟漪来反演行星内部结构，提出者 Mankovich 非常年轻，现在还是博士后，这种方法借鉴自星震学，可见行星和恒星研究之间是存在桥梁的；此外，还可以对木星进行持续观测，捕捉其表面发生的某些天气现象，比如周期性的对流爆发（这一研究已经能够模拟出来了，但是仍需要进一步研究）。}
\SourceListItem{· 天王星与海王星}
\SourceParagraph{对冰巨星的研究基本还在用冷战期间的探测器飞掠数据和望远镜观测，我们对其知之甚少。目前可以知道的是其质量较小，而核心质量比较大，表面有云和风暴，与木星和土星相比缺少氨气等等，远没有详细的温度和成分的垂直剖面可以提供，我们对其大气环流也知之甚少。但是也有一定的研究。}
\SourceParagraph{计划上，天问四号原打算分离一部分到海王星，成为第一个海王星环绕器的，后来这一计划由于风险较大、完成周期长，加上行星大气缺乏学科带头人、行星大气声量不如行星地质，所以该方案被取消，改为了对木卫四的着陆探测（其实研究木卫四的科学价值不如木卫二和木卫三，只是 NASA 和 ESA 已经有探测器去木卫二和木卫三了，因此，为形成差异化探测目标，任务转向了木卫四。）}
\SourceListItem{· 冥王星}
\SourceParagraph{对冥王星的近距离探测来自 2015 年新视野号（New Horizons）的飞掠。冥王星的大气其实不算稀薄，主要成分为氮气、甲烷、一氧化碳，目前对其大气的研究主要关注的是化学成分，也就是氮气和甲烷在紫外线照射下如何发生光化学反应，形成光化学烟雾（蓝色雾霾）。}
\SourceParagraph{系内行星及卫星大气的研究还是非常丰富的，未来十年至二十年的大量探测任务也将会继续推进我们对其的理解。当然，行星大气研究并不局限于太阳系内行星，与此同时，近年来对系外行星的研究也在井喷。系外行星直到 20 世纪 90 年代才首次被发现，后来逐渐增多，目前已有数千颗系外行星被发现，甚至有推测认为大多数恒星其实都带有系外行星。对系外行星的发现主要依赖于凌星法和径向速度法，除此之外也有直接成像法和引力透镜法，然而目前的技术直接成像困难很大，不过，据悉相关技术问题已取得突破，南大参与的“觅音计划”就是希望能直接成像探测系外宜居行星，目前实验星已经准备完毕，计划近期发射，后续数年陆续发射组件阵列望远镜后才正式开始观测。系外行星种类颇多，可以大体分为三类：热木星、亚海王星和类地行星。热木星因为质量大、轨道半径小（因而周期很小），所以观测起来非常方便，目前对热木星的研究包括但不限于成分、大气环流、大气逃逸等，由于数据多、验证容易，因此非常热门；亚海王星即质量在地球和海王星之间的行星，我对其了解不多；类地行星则不必多说，目前对类地行星的观测主要集中于大气成分、温度等，基于少量的观测约束来进行建模，估算其宜居程度。目前有不少研究认为传统意义上的“宜居带”可以根据行星的轨道或者成分等因素造成的气候效应来进行拓展，这一点气候/大气辐射/气溶胶等领域的研究者可以关注。}
\SourceParagraph{除了上述这些直接跟大气化学、大气物理、大气动力学、气候学相关的研究外，一些按理说属于大气科学范畴、但在南大大气尚未充分发展的方向也在蓬勃发展。比如高层大气或者空间天气，即研究电离层、热层及磁层的方向。对高层大气的研究基本上是空间科学相关学科的人在做（甚至就在隔壁天空院或者深空院），这方面比较出名的是科大。但是其研究方法主要也是基于流体力学的数值模拟，只是需要加入洛伦兹力和麦克斯韦方程组（MHD）。对地球及其他行星的高层大气进行探测也是一些探测任务的重点，甚至有些探测器并非以行星探测为主要目标也可以用来探测行星，比如某些太阳探测器（如羲和号）在借助巨行星（主要是木星）进行引力助推时也能顺便研究木星的磁层。对高层大气的研究主要有磁重联、磁层与太阳风相互作用（会形成一些波动）、大气逃逸（中性逃逸、离子逃逸等）等，相关的研究也很多。}
\SourceParagraph{如果要从 GAFD 的角度来看行星，那么相邻的方向就更加广阔了，研究行星气候的长期演化（如古火星/金星气候）、研究卫星（泰坦的甲烷大气、以及冰卫星的冰下海洋）、研究行星形成过程中的流体力学问题（原行星盘、行星大气吸积）、研究行星内部发电机（金属氢/铁核自转形成磁场）、研究恒星大气（自带热源、没有外来热源的流体，通常需要耦合磁场，因此恒星大气会比地球的复杂不少，比如旋转 Rayleigh-Bernard 对流、开尔文波等等，某些褐矮星其实可以结合类木行星与恒星建模研究），甚至是地质相关的行星演化（个人拙见，lava world 与岩浆洋的流体动力学应该有相似之处，只不过一个是长期的，一个只是行星形成的一个阶段）与地幔对流（虽然是固体物理，但是都是连续介质对流，其实不无相似之处）。}
\SourceParagraph{对这一庞大领域的介绍暂且到这里，接下来该讲讲国内外行星大气的情况了。国内研究行星大气主要有如下机构：北京大学物理学院大气与海洋科学系，主要做系外行星大气，偏气候，不过几位老师的视野都很广，也有做木星等（过去 Showman 在这做过讲座教授，带过一些做木星的学者）；澳门科技大学，依托月球与行星科学国家重点实验室，澳科大在行星科学方面实力较强，尤其是月球、火星等方面；南京大学，前沿科学学院深空探测科学与技术研究院，是与澳门科技大学共建月球与行星科学国家重点实验室的单位，研究领域涵盖火星、木星、冰卫星、小行星、原行星盘等等，天文与空间科学学院对行星大气的研究主要是偏系外行星观测的；中科院上海天文台，主要做木星大气、系外宜居行星探测、高层大气等；中科院大气物理研究所，我了解的主要是做火星大气建模的；上海交通大学李政道研究所有做系外行星大气环流建模的老师；关于原行星盘，做的主要是天文那边的老师，比如中山大学、中科院紫金山天文台、中科院云南天文台、北京师范大学、台湾中研院等；关于高层大气及大气逃逸，主要是空间科学那边的老师，比如中国科学技术大学、北京航空航天大学等。}
\SourceParagraph{国外的研究就更多了，规模比较大的有 UCSC、Exeter 大学等，研究学者则遍布 MIT、Harvard、Caltech、UMich、Oxford 等，甚至欧洲大陆上也有不少研究所从事相关研究，比如 Laboratoire de météorologie dynamique（LMD）、马普所等。}
\SourceParagraph{具体到研究者的话，国内研究系外行星大气首推北大，目前主力的杨军老师和丁峰老师实力都非常强悍，Daniel Koll 的流体动力学也很强，三位老师各有特点，兴趣广泛；上交李所的谭先瑜老师主要做系外行星大气环流建模，成果也很好，他的博士后赖燕红也是做系外行星大气环流的，赖东老师有做地外行星动力学和原行星盘，上交天文本身也有一位外籍学者从事系外行星研究，我了解不多，大家可以自行了解；上台的孔大力台长主要做木星内部及大气，博士后连雨辰做木星建模和系外行星大气建模，黄辰亮老师主要做系外大气逃逸，现在也在负责一些卫星相关的工程事务，袁靓亮老师主要做行星高层大气与电离层；南大天空院的刘慧根老师主要做系外行星大气探测，谢基伟老师则主要做系外行星普查，深空院的倪冬冬老师研究方向很多，气态行星内部和大气、冰卫星、发电机理论等都有，余亮亮老师有做火星大气，崔灿老师主要做原行星盘，大气科学学院的傅豪老师虽然本身是做热带动力学的，但是对行星大气很感兴趣，也在推动大气科学学院发展这一方向，丁家琛老师有做过火星大气光学相关的研究；澳门科技大学的蔡涛老师主要做木星大气，目前也在做木星极涡，肖静老师主要做火星大气建模。以上老师有很多是南大出身的，甚至有几位是南大大气出身的。}
\SourceParagraph{海外就更丰富了，Caltech 的 Andrew Ingersoll 是行星大气领域的泰斗，虽然年事已高，但是仍然在研究一线，许多在美国做行星大气研究的老师都是他的学生；MIT 的康婉莹是新生代老师中做地外行星的代表性学者，冰卫星、木星大气、系外行星都做；UMich 的 李 成 则 是 气 态 巨 行 星 领 域 的 代 表 性 学 者 ， 领 衔 Juno 任 务 ； Harvard 的 Robin Wordsworth 主要做系外行星气候，他们那的 Eli Tziperman 也是一位兴趣广泛的资深学者；以色列 Weizmann 研究所的 Yohai Kaspi 做木星大气动力学、系外行星大气，也做地球上的 GFD，他的组很大，实力较强；UCB 的 Philip Marcus 主要做流体力学，但是也有很多气态行星相关的动力学研究；UCLA 的曹浩主要做木星内部的 MHD 模拟，也做木星大气；UCSD 的 Lia Siegelman 主要从事海洋研究，但是也有很多木星大气的研究；LMD 的 François Forget 主要做类地行星的大气；Oxford 的 Peter Read 也做气态行星大气，不过他也到了退休年纪了；Geoffrey Vallis 目前在 Exeter，他本是做大气环流的，近期对热带对流、行星大气环流都很感兴趣，不过年龄也比较大了；University of Houston 的 Liming Li 是南大大气校友，也是做木星的华人里资历比较老的，经历过 Galileo、Cassini、Juno 等任务……除此之外还有很多人活跃在这一领域，恕我不能一一列出。}
\SourceParagraph{做行星大气最大的问题是很多理论没有前人的完整理论体系，没法验证，也很难说短期内有什么应用价值，但是因此受到的质疑和挑战相对较少。这也意味着如果离开了相关探测任务，经费会非常少，尤其是在美国。所以很多做行星大气，尤其是系外行星大气的老师，除了做理论和数值模拟的探索性工作之外，也会兼顾其他较容易获得经费支持的方向，比如环境、气象或者跟探测任务直接相关的方向，如火星、金星、木星等的大气。}
\SourceParagraph{对想要进入这一领域的同学来说（P.S. 由于我个人做的是动力学所以后面可能更偏行星大气动力学的路径介绍，如果想做化学或者辐射之类的可以单独联系我），首先需要保证基础概念的熟悉，比如如果做大气环流或者流体相关的，流体力学和动力气象基本的语言和方法论要比较熟悉。更前置的内容，比如力热光电和微积分这种，需要熟练运用，但是不需要奇技淫巧，所有不建议为了刷分而“过拟合”，去做很多偏题怪题，建议关注原理和基本概念/公式。理论上如果学有余力的话，这块可以加速，几个月应该就能解决，然后跳过概论，直接进入专业核心课学习（当然前提是你知道自己的兴趣在哪，概论课程理论上就是带你简要地走过各个方向，让你挑选感兴趣的方向），比如流体力学和动力气象。流体力学课本里面可能比较偏力学，这些可以不用太在意，主要是熟悉语言和方法论，同样，做题在这里仅供熟悉语言和帮助理解，如果跟着推一遍正文内容可以达到同样的效果，做题就不是必要的。动力气象前面基本上就是流体力学，后面是气象这边对 NS 方程的化简和分析，这个比较重要。如果看的是 Holton 那本书，最后几章比较偏诊断、数值模式和同化的看不看都行，毕竟行星里面做环流诊断的工作其实不多，数值模式和同化则有另外的书。然后还需要一些数值天气预报，虽然这个课程听起来是做地球预报的，但是本质是将方程离散化数值求解，所以对于做行星大气非常重要，即使想要做纯理论或者纯观测，最后也离不开跑模式来检验。再之后就能够看懂文献了，可以找老师要一些文献来读，然后也可以多听听学术报告，参加一下学术会议。刚开始可能感觉啥也听不懂，不过逐渐地你会发现很多东西会在报告和论文里不断重复，最后很多开始完全不知所云的东西对你来说不过是重复了，于是报告体验就会好很多。在这一过程中，如果主动思考报告的内容，然后跟老师交流交流，总能找到感兴趣的科学问题的，可以以科学问题为抓手开展研究，这样就算是进入了一个领域。所以简而言之，基本的路径都是基础课→专业课→文献/报告→开展研究，如果对专业课要看哪些书学哪些内容，文献要看哪些不太清楚，可以咨询老师，写邮件或者约时间聊天都可以。行星大气还没有非常体系化的专业课，所以还是得从其他领域的专业课入手，就比如我刚才说的大气动力学，或者大气辐射/恒星相关课程。做行星大气化学的话大气化学这门课用处不大，主要是地球上的大气化学研究的东西在其他星球都不多，臭氧、有机物这种都很少，这种的可能直接看文献，边看边学比较合适。当然做行星大气也要考虑好做观测、做模拟还是做偏理论，这直接影响你需要跟哪些老师。做观测的话基本上就是学习处理数据了，做理论的话就需要更多的物理，当然二者都需要涉及一些模拟，做模拟的话则要涉及二者（}
\SourceParagraph{未完待续…………}
